% Native hgraph authoring; original AimPaper statements are checked verbatim against AimPaper/main.tex.
% Source locators for the corrections tab; no Blueprint statements are retained.
% AimPaper: 6.1 theorem 1606 1621
% AimPaper: 6.12 theorem 1780 1799
\section{No-crossing criteria for filtered extensions}

% AimPaper: 2.6 corollary 327 330
\begin{corollary}[Vanishing below the Adams filtration]
\label{paper:2.6}
\source{AimPaper Corollary 2.6}
\level{coarse}
From \cite{LWX126}, Corollary 2.6.\par
% BEGIN AIMPAPER STATEMENT
If the Adams filtration of $f$ is $k$, then
    $d^f_i=0$ for $i<k$.
% END AIMPAPER STATEMENT
\uses{paper:2.5,sem:AdamsSynthetic:filtration-comparison}
% Blueprint note: Mathematical dependency review: [AimPaper source, lines 332-339](paper/main.tex.html#L332). Factor f into positive-filtration maps, then apply Proposition 2.5(1) to rule out shorter extension differentials.



% Correspondence: fixed original representative-form statement proved.
% Blueprint note: Original statement and context: paper/main.tex.html#L327; paper/2412.10879.pdf.
% Blueprint note: Original-statement interface: \texttt{KIP126.ProofTools.Paper.adamsFiltrationVanishing}. The public statement is proved without changing its type or premises.
% Blueprint note: The proof lifts the map along the actual target Adams tower and proves that filtration adds under composition. Every shorter extension target has zero associated grade.
\lean{KIP126.ProofTools.Paper.adamsFiltrationVanishing}
\end{corollary}
\begin{proof}
\textit{(See \cite{LWX126}, Corollary 2.6.)} \href{paper/main.tex.html#L332}{Original proof}.

When $k>0$, it is clear that $d^f_0=0$, which means that for any $s$ and $[x]\in F_s\pi_*(X)$, we have $f([x])\in F_{s+1}\pi_*(Y)$. Since $\AF(f)=k$, we know that there exist $k$ maps
    $$X=X_0\fto{f_1}X_1\fto{f_2}\cdots\fto{f_k}X_k=Y$$
    such that each $f_i$ induces trivial map in $\HF$-homology (hence $\AF(f_i)>0$) and $f$ is homotopic to the composition $f_k\circ \cdots\circ f_1.$
    Then for any $s$ and $[x]\in F_s\pi_*(X)$, inductively we have
    $$f_1([x])\in F_{s+1}\pi_*(X_1)$$
    $$f_2\circ f_1([x])\in F_{s+2}\pi_*(X_2)$$
    $$\cdots$$
    $$f([x])=f_k\circ\cdots\circ f_1([x])\in F_{s+k}\pi_*(X_k)=F_{s+k}\pi_*(Y).$$
    By Proposition \ref{prop:8154e6f1}.(1) this means that $d_i^fx=0$ for any $x\in E_\infty^{*,*}(X)$ and $i<k$.
\end{proof}

% AimPaper: 2.7 definition 344 352
\begin{definition}[Crossings of extensions]
\label{paper:2.7}
\source{AimPaper Definition 2.7}
\level{coarse}
From \cite{LWX126}, Definition 2.7.\par
% BEGIN AIMPAPER STATEMENT
\label{def:98skj23}
\begin{enumerate}[leftmargin=*]
    \item For $x\in E_\infty^{s,t}(X)$, we say that the an $f$-extension $d_n^{f}(x)=y$ has a crossing that hits Adams filtration $p$ for some $p \leq s+n$, if there exists an element $x^\prime\in E_\infty^{s+a,t+a}(X)$ with $a>0$ and an $f$-extension $d_m^{f}(x^\prime)=y^\prime$ for $0\neq y^\prime\in E^{s+a+m, t+a+m}_\infty(Y)$ such that
    $$p = \textup{AF}(y^\prime) = s+a+m\le \textup{AF}(y) = s+n.$$
    \item We say that an $f$-extension $d_n^{f}(x)=y$ has no crossing that hits the range AF $\ge p$ if there does not exist an element $x^\prime\in E_\infty^{s+a,t+a}(X)$ with $a>0$ and an $f$-extension $d_m^{f}(x^\prime)=y^\prime$ for $0\neq y^\prime\in E^{s+a+m, t+a+m}_\infty(Y)$ such that
    $$p\le \textup{AF}(y^\prime) = s+a+m\le \textup{AF}(y) = s+n.$$
    \item We say that an $f$-extension $d_n^{f}(x)=y$ has no crossing if it has no crossing that hits the range AF$\ge \AF(x) + 1$.
\end{enumerate}
% END AIMPAPER STATEMENT
\uses{paper:2.3}
% Blueprint note: Mathematical dependency review: [AimPaper source, lines 344-352](paper/main.tex.html#L344). A crossing is a competing extension from a higher-filtration source, using the extension relation already defined.

% Blueprint note: The generic filtered-map crossing predicate belongs to \ref{sem:SpectralSequence:crossing-transport}. The bindings here are the two stated no-crossing predicates with an arbitrary filtration threshold.


% Correspondence: partial.
% Blueprint note: Public representative-level definitions \texttt{FilteredNoCrossingRange} and \texttt{ExtensionNoCrossingRange} allow arbitrary filtration thresholds. Their comparison with \texttt{DifferentialDatum} and \texttt{HasCrossingAt} for a general classical $f$-extension spectral sequence has not yet been assembled into a complete public bridge, so this correspondence remains partial.
% Blueprint note: Direct formal counterparts: \texttt{KIP126.ProofTools.Paper.FilteredNoCrossingRange}, \texttt{KIP126.ProofTools.Paper.ExtensionNoCrossingRange}.
\lean{KIP126.ProofTools.Paper.FilteredNoCrossingRange,KIP126.ProofTools.Paper.ExtensionNoCrossingRange}

\sketch
% Blueprint note: Original statement and context: paper/main.tex.html#L344; paper/2412.10879.pdf.
\end{definition}

% AimPaper: 2.10 proposition 408 416
\begin{proposition}[Characterization of extensions without crossings]
\label{paper:2.10}
\source{AimPaper Proposition 2.10}
\level{coarse}
From \cite{LWX126}, Proposition 2.10.\par
% BEGIN AIMPAPER STATEMENT
\label{prop:i8r47oe}
    An $f$-extension from $x$ to $y$ has no crossing in the range $\AF\ge p$ if and only if for all $[x]\in \{x\}$ such that $\AF(f[x])\ge p$ we have $f[x]\in \{y\}$. In particular, we have
    \begin{enumerate}
        \item An $f$-extension from $x$ to $y$ has no crossing if and only if for all $[x]\in \{x\}$ we have $f[x]\in \{y\}$.
        \item An $f$-extension $d_n^{f}(x)=0$ has no crossing if and only if for all $[x]\in \{x\}$,
        $$\AF(f[x]) >\AF(x)+n.$$
        \item If $\AF(f)=n$, then all $d^f_n$-differentials have no crossing in the $f$-ESS.
    \end{enumerate}
% END AIMPAPER STATEMENT
\uses{paper:2.5,paper:2.7}
% Blueprint note: Mathematical dependency review: [AimPaper source, lines 418-430](paper/main.tex.html#L418). Proposition 2.5(3) and the definition of a crossing turn universal representative behavior into the no-crossing criterion.



% Correspondence: public counterpart, within the scope described below.
% Blueprint note: The full representative-level equivalence and its three clauses are public for arbitrary filtered abelian groups. An $E_\infty$ formulation for the same synthetic map is also available. This does not establish the missing general classical $f$-extension spectral sequence construction in Definition~2.1.
% Blueprint note: Formal counterparts: \item \texttt{KIP126.ProofTools.Paper.filtered\_extension\_noCrossingRange\_iff} \item \texttt{KIP126.ProofTools.Paper.filtered\_extension\_noCrossing\_iff} \item \texttt{KIP126.ProofTools.Paper.filtered\_zero\_extension\_noCrossing\_iff} \item \texttt{KIP126.ProofTools.Paper.filtered\_extension\_noCrossing\_of\_filtration\_bound}
\lean{KIP126.ProofTools.Paper.filtered_extension_noCrossingRange_iff,KIP126.ProofTools.Paper.filtered_extension_noCrossing_iff,KIP126.ProofTools.Paper.filtered_zero_extension_noCrossing_iff,KIP126.ProofTools.Paper.filtered_extension_noCrossing_of_filtration_bound}

% Blueprint note: Original statement and context: paper/main.tex.html#L408; paper/2412.10879.pdf.
\end{proposition}
\begin{proof}
\textit{(See \cite{LWX126}, Proposition 2.10.)} \href{paper/main.tex.html#L418}{Original proof}.

First we prove the only if part. Assume by contradiction that there exists $[x]\in \{x\}$ such that $f[x]=[y']\notin \{y\}$ is detected by $y'\in E_\infty(Y)$ with $\AF(y')\ge p$. Then there is an $f$ extension from $x$ to $y'$. By Proposition \ref{prop:8154e6f1}~(3), $y'$ (if $\AF(y)>\AF(y')$)  or $y-y'$ (if $\AF(y)=\AF(y')$) or $y$ (if $\AF(y)<\AF(y')$) should be hit by a shorter $d^f$ differential. This shorter differential is a crossing of the $f$-extension from $x$ to $y$ that hits $\AF\ge p$, which is a contradiction. Now we prove the if part. Assume that the $f$-extension has a crossing from $x'$ to $y'$ with $p\le \AF(y')\le \AF(y)$. Then there exists $[x']\in \{x'\}$ such that $f([x'])\in \{y'\}$. Note that $[x]+[x']\in \{x\}$ since $\AF(x')>\AF(x)$, we have
    $$f([x]+[x'])=f([x])+f([x'])=[y]+[y']\in\begin{cases}
        \{y'\} & \text{ if } \AF(y')<\AF(y)\\
        \{y+y'\} & \text{ if } \AF(y')=\AF(y).\\
    \end{cases}$$
    However, by assumption we should have $f([x]+[x'])\in \{y\}$ and it is contradictory to the equation above in either case.

    Hence the main statement is true.  Part (1) (2) are special cases of the main statement and Part (3) holds for degree reasons.
\end{proof}

\section{Commutative squares and extension comparison}

% AimPaper: 2.12 theorem 461 481
\begin{theorem}[The extension theorem for a commutative square]
\label{paper:2.12}
\source{AimPaper Theorem 2.12}
\level{coarse}
From \cite{LWX126}, Theorem 2.12.\par
% BEGIN AIMPAPER STATEMENT
\label{thm:4114f70c}
    Consider a homotopy commutative diagram of spectra
    $$\xymatrix{
    X \ar[r]^{f}\ar[d]_{p} & Y\ar[d]^{q}\\
    Z \ar[r]_{g} & W
    }$$
    Suppose $m,n,l\ge 0$, $0<k\le m+l-n$,
    $$\begin{array}{ll}
        x\in E^{s,t}_\infty(X) & y\in E^{s+n,t+n}_\infty(Y),\\
        z\in E^{s+m,t+m}_\infty(Z) & w\in E^{s+m+l,t+m+l}_\infty(W),
    \end{array}$$
    and
    \begin{enumerate}
        \item $d^{f}_{n}(x)=y$,
        \item $d^{p}_{m}(x)=z$,
        \item the differential in (1) or the differential in (2) has no crossing,
        \item $d^{g}_{l}(z)=w$, and has no crossing that hits the range AF $\ge s+n+k$,
        \item $d^{q}_{k-1} y=0$, and has no crossing.
    \end{enumerate}
    Then $d^{q}_{m+l-n}(y)=w$. (See Figure \ref{fig:4f154a8}.)
% END AIMPAPER STATEMENT
\uses{paper:2.10}
% Blueprint note: Mathematical dependency review: [AimPaper source, lines 483-505](paper/main.tex.html#L483). Use the representative criterion to choose a common source and transport its images around the commutative square.



% Correspondence: fixed original representative-form statement proved.
% Blueprint note: Original statement and context: paper/main.tex.html#L461; paper/2412.10879.pdf.
% Blueprint note: Original-statement interface: \texttt{KIP126.ProofTools.Paper.extensionSquare}. The public statement is proved without changing its type or premises.
% Blueprint note: The proved finite-filtration representative criterion synchronizes the two source edges. The zero $q$-extension puts the same representative above $s+n+k$; only the stated threshold no-crossing condition for $g$ is then used.
\lean{KIP126.ProofTools.Paper.extensionSquare}
\end{theorem}
\begin{proof}
\textit{(See \cite{LWX126}, Theorem 2.12.)} \href{paper/main.tex.html#L483}{Original proof}.

First we find a representative $[x]$ of $x$ such that
    \begin{equation}\label{eq:9d45feac}
        p[x]\in \{z\}
    \end{equation}
    and
    \begin{equation}\label{eq:36e72616}
        f[x]\in \{y\}.
    \end{equation}
    If (1) has no crossing, by (2) we can pick $[x]\in \{x\}$ such that (\ref{eq:9d45feac}) holds;
    If (2) has no crossing, by (1) we can pick $[x]\in \{x\}$ such that  (\ref{eq:36e72616}) holds.
    In both cases the other equation also holds because of the no crossing condition.

    By (5), (\ref{eq:36e72616}) and Proposition~\ref{prop:i8r47oe}~(2) we know that $qf[x]$ has AF $\ge s+n+k$. Since the diagram on the Adams $E_\infty$-pages commutes, we have
    $$AF(gp[x])=AF(qf[x])\ge s+n+k.$$
    Combining with (4), (\ref{eq:9d45feac}) and Proposition~\ref{prop:i8r47oe} we have
    $$gp[x]\in \{w\}.$$
    Therefore, $qf[x]\in \{w\}$ and by (\ref{eq:36e72616}) there is a $q$-extension from $y$ to $w$.
\end{proof}

% AimPaper: 2.15 corollary 555 574
\begin{corollary}[A square of extensions without crossings]
\label{paper:2.15}
\source{AimPaper Corollary 2.15}
\level{coarse}
From \cite{LWX126}, Corollary 2.15.\par
% BEGIN AIMPAPER STATEMENT
\label{cor:0012nik}
    Consider a homotopy commutative diagram of spectra
    $$\xymatrix{
    X \ar[r]^{f}\ar[d]_{p} & Y\ar[d]^{q}\\
    Z \ar[r]_{g} & W
    }$$
    Suppose $m,n,l\ge 0$, 
    $$\begin{array}{ll}
        x\in E^{s,t}_\infty(X) & y\in E^{s+n,t+n}_\infty(Y),\\
        z\in E^{s+m,t+m}_\infty(Z) & w\in E^{s+m+l,t+m+l}_\infty(W),
    \end{array}$$
    and
    \begin{enumerate}
        \item $d^{f}_{n}(x)=y$,
        \item $d^{p}_{m}(x)=z$,
        \item the differential in (1) or the differential in (2) has no crossing,
        \item $d^{g}_{l}(z)=w$, and has no crossing.
    \end{enumerate}
    Then $d^{q}_{m+l-n}(y)=w$.
% END AIMPAPER STATEMENT
\uses{paper:2.12}
% Blueprint note: Mathematical dependency review: [AimPaper source, lines 555-574](paper/main.tex.html#L555). Specialize the square theorem to the full no-crossing hypothesis on the lower extension.



% Correspondence: fixed original representative-form statement proved.
% Blueprint note: Original statement and context: paper/main.tex.html#L555; paper/2412.10879.pdf.
% Blueprint note: Original-statement interface: \texttt{KIP126.ProofTools.Paper.noCrossingExtensionSquare}. The public statement is proved without changing its type or premises.
% Blueprint note: The proved filtered-group no-crossing criterion is specialized directly to the actual classical Adams tower. It synchronizes a common source representative and the lower target, and the same commutative square gives the right extension. No synthetic-to-classical comparison or additional premise is needed.
\lean{KIP126.ProofTools.Paper.noCrossingExtensionSquare}
\end{corollary}
\begin{proof}
\textit{(See \cite{LWX126}, Corollary 2.15.)} \href{paper/main.tex.html#L555}{Original statement and context}.

Choose one representative $[x]$ for which both $f[x]$ and $p[x]$ have the prescribed detectors $y$ and $z$. Such a common choice exists because one of these two extensions has no crossing: choose a representative for the other extension, and apply Proposition~\ref{paper:2.10} to the crossing-free one.

The crossing-free $g$-extension then implies that $gp[x]$ is detected by $w$. Homotopy commutativity gives $qf[x]=gp[x]$. Thus $f[x]$, detected by $y$, has $q$-image detected by $w$, which is exactly the required extension. Its length is the filtration difference $(s+m+l)-(s+n)=m+l-n$.
\end{proof}

% AimPaper: 2.16 corollary 576 591
\begin{corollary}[Extensions in a commutative triangle]
\label{paper:2.16}
\source{AimPaper Corollary 2.16}
\level{coarse}
From \cite{LWX126}, Corollary 2.16.\par
% BEGIN AIMPAPER STATEMENT
\label{cor:166dc180}
    Consider a homotopy commutative diagram of spectra
    $$\xymatrix{
    X \ar[r]^{f}\ar[dr]_{p} & Y\ar[d]^{q}\\
    & Z
    }$$
    Suppose $n,m\ge 0$, 
    $$x\in E^{s,t}_\infty(X), ~y\in E^{s+n,t+n}_\infty(Y), ~z\in E^{s+m,t+m}_\infty(Z)$$
    and
    \begin{enumerate}
        \item $d^{f}_{n}(x)=y$,
        \item $d^{p}_{m}(x)=z$,
        \item the differential in (1) or the differential in (2) has no crossing.
    \end{enumerate}
    Then $d^{q}_{m-n}(y)=z$.
% END AIMPAPER STATEMENT
\uses{paper:2.15}
% Blueprint note: Mathematical dependency review: [AimPaper source, lines 593-595](paper/main.tex.html#L593). Specialize Corollary 2.15 to an identity lower horizontal map.



% Correspondence: fixed original representative-form statement proved.
% Blueprint note: Original statement and context: paper/main.tex.html#L576; paper/2412.10879.pdf.
% Blueprint note: Original-statement interface: \texttt{KIP126.ProofTools.Paper.extensionTriangle}. The public statement is proved without changing its type or premises.
% Blueprint note: The common-source representative criterion proves the third extension directly. The conclusion keeps length $m-n$, including the signed-length cases of the fixed statement.
\lean{KIP126.ProofTools.Paper.extensionTriangle}
\end{corollary}
\begin{proof}
\textit{(See \cite{LWX126}, Corollary 2.16.)} \href{paper/main.tex.html#L593}{Original proof}.

This is a special case of Corollary \ref{cor:0012nik} when $Z=W$ and $g=id_Z$.
\end{proof}

% AimPaper: 2.17 corollary 597 611
\begin{corollary}[Composition of extensions]
\label{paper:2.17}
\source{AimPaper Corollary 2.17}
\level{coarse}
From \cite{LWX126}, Corollary 2.17.\par
% BEGIN AIMPAPER STATEMENT
\label{cor:290d35ce}
    Consider a homotopy commutative diagram of spectra
    $$\xymatrix{
    X \ar[d]_{p}\ar[dr]^{q} & \\
    Z\ar[r]_{g} & W
    }$$
    Suppose $n,m\ge 0$, 
    $$x\in E^{s,t}_\infty(X), ~z\in E^{s+m,t+m}_\infty(Y), ~w\in E^{s+m+l,t+m+l}_\infty(W)$$
    and
    \begin{enumerate}
        \item $d^{p}_{m}(x)=z$,
        \item $d^{g}_{l}(z)=w$, and has no crossing.
    \end{enumerate}
    Then $d^{q}_{m+l}(x)=w$.
% END AIMPAPER STATEMENT
\uses{paper:2.15}
% Blueprint note: Mathematical dependency review: [AimPaper source, lines 613-615](paper/main.tex.html#L613). Specialize Corollary 2.15 to an identity upper horizontal map.



% Correspondence: fixed original representative-form statement proved.
% Blueprint note: Original statement and context: paper/main.tex.html#L597; paper/2412.10879.pdf.
% Blueprint note: Original-statement interface: \texttt{KIP126.ProofTools.Paper.extensionComposition}. The public statement is proved without changing its type or premises.
% Blueprint note: The printed occurrence of $Y$ for $z$ is read as $Z$, as forced by the diagram; the unused $n$ and the printed $n,m\ge0$ are retained without adding $l\ge0$. See the statement-typing correction record.
% Blueprint note: Apply the lower no-crossing criterion to the actual image of a source representative, then use the composite equality. No additional condition $l\ge0$ is imposed.
\lean{KIP126.ProofTools.Paper.extensionComposition}
\end{corollary}
\begin{proof}
\textit{(See \cite{LWX126}, Corollary 2.17.)} \href{paper/main.tex.html#L613}{Original proof}.

This is a special case of Corollary \ref{cor:0012nik} when $X=Y$ and $f=id_X$.
\end{proof}

\section{Functoriality and exactness of extension spectral sequences}

% AimPaper: 2.18 corollary 617 626
\begin{corollary}[Morphisms of extension spectral sequences]
\label{paper:2.18}
\source{AimPaper Corollary 2.18}
\level{coarse}
From \cite{LWX126}, Corollary 2.18.\par
% BEGIN AIMPAPER STATEMENT
\label{cor:e7b20ae2}
    Consider a homotopy commutative diagram of spectra
    $$\xymatrix{
    X \ar[r]^{f}\ar[d]_{p} & Y\ar[d]^{q}\\
    Z \ar[r]_{g} & W
    }.$$
    Assume $\ESS{p}_0^{*,*}=\ESS{p}_r^{*,*}$ and $\ESS{q}_0^{*,*}=\ESS{q}_r^{*,*}$ for some $r\ge 0$. Then 
    $$(d^p_r, d^q_r): E_\infty^{*,*}(X)\oplus E_\infty^{*,*}(Y)\to E_\infty^{*,*}(Z)\oplus E_\infty^{*,*}(W)$$
    induces a map from the $f$-ESS to the $g$-ESS.
% END AIMPAPER STATEMENT
\uses{paper:2.15,paper:2.1}
% Blueprint note: Mathematical dependency review: [AimPaper source, lines 628-630](paper/main.tex.html#L628). Apply the square corollary inductively to obtain commuting maps on every page of the extension spectral sequences.



% Correspondence: fixed original representative-form statement proved.
% Blueprint note: Original statement and context: paper/main.tex.html#L617; paper/2412.10879.pdf.
% Blueprint note: Original-statement interface: \texttt{KIP126.ProofTools.Paper.extensionSpectralSequenceMorphism}. The public statement is proved without changing its type or premises.
% Blueprint note: Finite induction on the preceding vanishing extension relations proves a filtration shift by $r$. The unique degree-$r$ extension values form additive maps whose exact graphs are the $p$- and $q$-relations; the same homotopy representative and square preserve every $f$-extension relation. This proves the fixed statement in its full representative presentation; the separate construction of a classical ESS remains scoped in Definition 2.1.
\lean{KIP126.ProofTools.Paper.extensionSpectralSequenceMorphism}
\end{corollary}
\begin{proof}
\textit{(See \cite{LWX126}, Corollary 2.18.)} \href{paper/main.tex.html#L628}{Original proof}.

Since $\ESS{p}_0^{*,*}=\ESS{p}_r^{*,*}$ and $\ESS{q}_0^{*,*}=\ESS{q}_r^{*,*}$, we know that $d^p_r$, $d^q_r$ have no crossings. By Corollary \ref{cor:0012nik}, we have $d^g_0\circ d^p_r = d^q_r\circ d^f_0$ and therefore $(d^p_r, d^q_r)$ induces a map from $\ESS{f}_1^{*,*}$ to $\ESS{g}_1^{*,*}$. Inductively we can apply Corollary \ref{cor:0012nik} again and show that $(d^p_r, d^q_r)$ induces a map from $\ESS{f}_n^{*,*}$ to $\ESS{g}_n^{*,*}$ and $d^g_n\circ d^p_r = d^q_r\circ d^f_n$ for all $n$.
\end{proof}

% AimPaper: 2.19 corollary 632 635
\begin{corollary}[Permanent cycles from a null composite]
\label{paper:2.19}
\source{AimPaper Corollary 2.19}
\level{coarse}
From \cite{LWX126}, Corollary 2.19.\par
% BEGIN AIMPAPER STATEMENT
\label{cor:aed3d1a4}
    If the composition $g\circ f$ of two maps $f:X\to Y$, $g:Y\to Z$ is trivial, and $d_n^f(x)=y$ for $x\in E_\infty^{s,t}(X)$ and $y\in E_\infty^{s+n,t+n}(Y)$, then $y$ is a permanent cycle in the $g$-ESS, i.e., we have
    $d^g_m(y)=0$ for all $m\ge 0$.
% END AIMPAPER STATEMENT
\uses{paper:2.15,paper:2.5}
% Blueprint note: Mathematical dependency review: [AimPaper source, lines 637-646](paper/main.tex.html#L637). The null-composite square proves permanence; the article also gives a representative proof using Proposition 2.5.



% Correspondence: fixed original representative-form statement proved.
% Blueprint note: Original statement and context: paper/main.tex.html#L632; paper/2412.10879.pdf.
% Blueprint note: Original-statement interface: \texttt{KIP126.ProofTools.Paper.permanentCycle\_of\_nullComposite}. The public statement is proved without changing its type or premises.
% Blueprint note: The original extension supplies a single representative in the actual image of $f$. Its composite with $g$ is zero, so that same representative witnesses every zero outgoing extension.
\lean{KIP126.ProofTools.Paper.permanentCycle_of_nullComposite}
\end{corollary}
\begin{proof}
\textit{(See \cite{LWX126}, Corollary 2.19.)} \href{paper/main.tex.html#L637}{Original proof}.

By Proposition~\ref{paper:2.5}, choose a representative $[x]$ for which $[y]=f[x]$ is detected by $y$. The null composite gives $g[y]=gf[x]=0$. This same representative therefore has zero image in every filtration quotient. By the representative definition of the extension spectral sequence, every outgoing $g$-extension differential on $y$ vanishes.
\end{proof}

% AimPaper: 2.20 proposition 648 654
\begin{proposition}[Exactness and extension boundaries]
\label{paper:2.20}
\source{AimPaper Proposition 2.20}
\level{coarse}
From \cite{LWX126}, Proposition 2.20.\par
% BEGIN AIMPAPER STATEMENT
\label{prop:cfe810af}
    Suppose that the sequence of spectra
    $$X\fto{f}Y\fto{g}Z$$
    induces a sequence of homotopy groups
    $$\pi_*X\fto{\pi_*f}\pi_*Y\fto{\pi_*g}\pi_*Z$$
    that is exact in $\pi_*Y$. Then all permanent cycles in $E^{*,*}_\infty(Y)$ of the $g$-ESS are boundaries in the $f$-ESS.
% END AIMPAPER STATEMENT
\uses{paper:2.1,paper:2.5,sem:SpectralSequence:completion,sem:AdamsSynthetic:selected-convergence}
% Blueprint note: Mathematical dependency review: [AimPaper source, lines 656-664](paper/main.tex.html#L656). Convergence of the filtered complex and exactness identify permanent cycles with extension boundaries; representatives give the alternative proof.



As throughout \cite{LWX126}, spectra are understood to be 2-completed,
connective and of finite type. The Lean binding expresses the same
conclusion through the representative criterion of
Proposition~\ref{paper:2.5}, with completeness of the middle homotopy
filtration, finite nonnegative filtration quotients and separation of the
target filtration explicit in its background inputs. Permanence still
permits a different representative on each page; zero labels are included.
The fixed Standard model's realization of this completed background and
its convergence hypotheses remains an upstream task, tracked in
issue 31.
% Blueprint note: The AimPaper statement is retained verbatim and directly binds its proved representative-form theorem. The explicit background inputs are retained; this binding does not assert their construction for the fixed Standard model.
% Blueprint note: The deferred completed-background migration is tracked in [issue 31](https://github.com/SII-MATH/Kervaire126Main/issues/31).
\lean{KIP126.ProofTools.Paper.extensionBoundary_of_exact}
\end{proposition}
\begin{proof}
\uses{paper:2.5,sem:SpectralSequence:completion,sem:AdamsSynthetic:selected-convergence}
\textit{(See \cite{LWX126}, Proposition 2.20.)} \href{paper/main.tex.html#L656}{Original proof}.

The representative argument makes the convergence step explicit. Fix
$y\in E_\infty^{s,d+s}(Y)$. For each integer $k$, let $S_k$ consist of the
residue classes in $\pi_dY/F^k\pi_dY$ represented by some $b$ representing
$y$ with $g_*b\in F^k\pi_dZ$. Permanence makes every $S_k$ nonempty.
The finite nonnegative quotients also make every integer-indexed quotient
finite, and the quotient transitions carry $S_l$ into $S_k$ for $k\leq l$.
The finite inverse-system principle gives a compatible family of such
residues. Completeness lifts it to one element $b\in\pi_dY$.

The residue at $k=s+1$ shows that $b$ still represents $y$. For every
$n\geq0$, its residue at $n$ agrees with that of a representative whose
image lies in $F^n\pi_dZ$. Filtration preservation then puts $g_*b$ in
$F^n\pi_dZ$ as well. Separation gives $g_*b=0$, and homotopy exactness
supplies $c\in\pi_dX$ with $f_*c=b$. At $a=\min(0,s)$ the source Adams
filtration is the whole homotopy group, so $c$ defines an
associated-graded class $x$. Proposition~\ref{paper:2.5} gives the required
$f$-extension from $x$ to $y$, so $y$ is an $f$-extension boundary.
\end{proof}

\section{Synthetic boundary maps and extension differentials}

% AimPaper: 4.2 proposition 1025 1029
\begin{proposition}[Extensions for $\lambda^n$ and $\rho$]
\label{paper:4.2}
\source{AimPaper Proposition 4.2}
\level{coarse}
From \cite{LWX126}, Proposition 4.2.\par
% BEGIN AIMPAPER STATEMENT
\label{prop:9770ae6e} 
The only nonzero differentials in the extension spectral sequences for the maps $\lambda^n$ and $\rho$ from Notation~\ref{not:deltaandrho} are the $d_0$'s:
$$d^{\lambda^n}_0=\lambda^n \ \text{and} \ d^{\rho}_0=\rho.$$ 
As a result, these $d_0$'s have no crossings.
% END AIMPAPER STATEMENT
\uses{paper:3.11,paper:3.12,paper:2.19}
% Blueprint note: Mathematical dependency review: [AimPaper source, lines 1030-1050](paper/main.tex.html#L1030). The infinite-page quotient formulas make the lambda/rho sequence exact; the null-composite result rules out higher extensions.



% Correspondence: public counterparts; scope is specified below.
% Blueprint note: The finite-quotient and infinite versions concern the same actual maps $\lambda$ and $\rho$: they exclude essential extensions of positive length, identify the actual $d_0$ equations, and establish the absence of crossings.
% Blueprint note: This entry is bound to the existing Lean declarations. Its status reports their local completion; the scope of correspondence with the article is specified above.
% Blueprint note: Bound declarations (private names are resolved through stable source references without changing their visibility):
% Blueprint note: \begin{itemize}
% Blueprint note: \item \texttt{KIP126.ProofTools.Paper.lambda\_rho\_extension\_collapse\_finite}
% Blueprint note: \item \texttt{KIP126.ProofTools.Paper.lambda\_rho\_extension\_collapse\_infinite}
% Blueprint note: \end{itemize}
\lean{KIP126.ProofTools.Paper.lambda_rho_extension_collapse_finite,KIP126.ProofTools.Paper.lambda_rho_extension_collapse_infinite}

% Blueprint note: The original article number is 4.2; other numbering is generated by the blueprint. Throughout the article, spectra are assumed to be 2-completed, connective, and of finite type. Read this statement together with its original context.
% Blueprint note: \href{paper/main.tex.html#L1025}{Original statement, hypotheses, and context} · \href{paper/2412.10879.pdf}{Article PDF}
\end{proposition}
\begin{proof}
\textit{(See \cite{LWX126}, Proposition 4.2.)} \href{paper/main.tex.html#L1030}{Original proof}.

Compute the three associated graded groups in the cofiber sequence using Propositions~\ref{paper:3.11} and~\ref{paper:3.12}. Write $k=t-w$. Outside $0\le k<m$ the middle group vanishes. For $0\le k<n$, the $\lambda^n$ source group vanishes and $\rho$ is the inclusion of the relevant cycle subquotients. For $n\le k<m$, the $\rho$ target group vanishes and $\lambda^n$ is the quotient map from a smaller boundary quotient onto the middle group. Thus the sequence of associated graded groups is exact in the middle in every degree.

The two composites are zero. Corollary~\ref{paper:2.19} makes the image of each preceding map a permanent cycle for the next extension spectral sequence. Since each nonzero graded $\lambda^n$ map is surjective and each nonzero graded $\rho$ map is injective in these ranges, no positive-length extension remains. Hence the only differentials are the indicated $d_0$ maps, which have no crossings. The untruncated case uses $Z_\infty$ in the same calculation.
\end{proof}

% AimPaper: 4.6 proposition 1068 1086
\begin{proposition}[Boundary extensions from Adams differentials]
\label{paper:4.6}
\source{AimPaper Proposition 4.6}
\level{coarse}
From \cite{LWX126}, Proposition 4.6.\par
% BEGIN AIMPAPER STATEMENT
\label{prop:6de7d130}
    Suppose in the classical Adams spectral sequence of $X$ we have $d_r(x)=y$, where $x\in Z_{r-1}^{s,t}(X)$ and $y\in Z_\infty^{s+r,t+r-1}(X)$. Consider the map
    $$\delta:\nu X/\lambda^n\to \Sus^{1,-n}\nu X/\lambda^{m-n}.$$
    \begin{enumerate}
        \item If $r\ge n+1$, then we view $x$ as an element of 
        $$E_\infty^{s,t,t}(\nu X/\lambda^n)\iso Z_n^{s,t}(X),$$
        and $\lambda^{r-n-1} y$ as an element of 
        $$E_\infty^{s+r,t+r-1, t+n}(\nu X/\lambda^{m-n})\iso Z_{m-r+1}^{s+r,t+r-1}(X)/B_{r-n}^{s+r,t+r-1}(X).$$
        We then have
        $$d^{\delta}_r(x)=\lambda^{r-n-1}y,$$
        which is trivial if $r>m$.
        \item If $r<n+1$, then we view $\lambda^{n+1-r}x$ as an element of 
        $$E_\infty^{s,t,t-n-1+r}(\nu X/\lambda^n)\iso Z^{s,t}_{r-1}(X)/B^{s,t}_{n+2-r}(X),$$
        and $y$ as an element of 
        $$E_\infty^{s+r,t+r-1, t+r-1}(\nu X)\iso Z_\infty^{s+r,t+r-1}(X).$$
        In this case, we have
        $$d^{\delta}_r(\lambda^{n+1-r}x)=y.$$
    \end{enumerate}
% END AIMPAPER STATEMENT
\uses{paper:3.8,paper:3.11,paper:3.12,paper:4.2,paper:2.16,paper:2.17}
% Blueprint note: Mathematical dependency review: [AimPaper source, lines 1087-1131](paper/main.tex.html#L1087). Start with the Bockstein comparison, induct on the quotient length using lambda and restriction triangles, and reduce the finite target by composition of extensions.



% Correspondence: public counterparts; scope is specified below.
% Blueprint note: The four public statements cover both ranges of differential lengths, for finite and absolute $\delta$-extensions.
% Blueprint note: This entry is bound to the existing Lean declarations. Its status reports their local completion; the scope of correspondence with the article is specified above.
% Blueprint note: Bound declarations (private names are resolved through stable source references without changing their visibility):
% Blueprint note: \begin{itemize}
% Blueprint note: \item \texttt{KIP126.ProofTools.Paper.syntheticDeltaExtension\_prop46\_above\_absolute}
% Blueprint note: \item \texttt{KIP126.ProofTools.Paper.syntheticDeltaExtension\_prop46\_below\_absolute}
% Blueprint note: \item \texttt{KIP126.ProofTools.Paper.syntheticDeltaExtension\_prop46\_above\_finite}
% Blueprint note: \item \texttt{KIP126.ProofTools.Paper.syntheticDeltaExtension\_prop46\_below\_finite}
% Blueprint note: \end{itemize}
\lean{KIP126.ProofTools.Paper.syntheticDeltaExtension_prop46_above_absolute,KIP126.ProofTools.Paper.syntheticDeltaExtension_prop46_below_absolute,KIP126.ProofTools.Paper.syntheticDeltaExtension_prop46_above_finite,KIP126.ProofTools.Paper.syntheticDeltaExtension_prop46_below_finite}

% Blueprint note: The original article number is 4.6; other numbering is generated by the blueprint. Throughout the article, spectra are assumed to be 2-completed, connective, and of finite type. Read this statement together with its original context.
% Blueprint note: \href{paper/main.tex.html#L1068}{Original statement, hypotheses, and context} · \href{paper/2412.10879.pdf}{Article PDF}
\end{proposition}
\begin{proof}
\textit{(See \cite{LWX126}, Proposition 4.6.)} \href{paper/main.tex.html#L1087}{Original proof}.

First take $m=\infty$ and write $\delta_n=\delta_{n,\infty}$. For $n=1$, Theorem~\ref{paper:3.8} identifies the connecting-map extension with the classical Adams differential. Induct on $n$, using the triangle comparing the quotients of lengths $n-1$ and $n$ and multiplication by $\lambda$.

For $r\le n$, the commutative-triangle extension rule transports the induction hypothesis to $\lambda^{n+1-r}x$ and gives target $y$. For $r\ge n+1$, the comparison gives $\lambda d_r^{\delta_n}(x)=\lambda^{r-n}y$ modulo classical $B_{r-1}$. The possible ambiguity in dividing this particular equality by $\lambda$ is contained in $B_{r-n+1}\subseteq B_{r-1}$, by the explicit infinity-page subquotients. This yields $d_r^{\delta_n}(x)=\lambda^{r-n-1}y$ in the required quotient, rather than an unrestricted cancellation theorem.

Induction on the differential length identifies $B_{r-1}$ here with the sum of the earlier $\delta_n$-extension images: each earlier classical boundary gives its corresponding boundary extension, and the same comparison gives the converse. Finally $\delta_{n,m}$ is $\delta_n$ followed by restriction to the finite target quotient. Proposition~\ref{paper:4.2} and the composition rule transport both formulas to that quotient. The target is zero when its $\lambda$-exponent reaches the quotient length, giving the stated vanishing range.
\end{proof}

% AimPaper: 4.10 corollary 1138 1144
\begin{corollary}[$\lambda$ multiples of boundary extensions]
\label{paper:4.10}
\source{AimPaper Corollary 4.10}
\level{coarse}
From \cite{LWX126}, Corollary 4.10.\par
% BEGIN AIMPAPER STATEMENT
\label{cor:2a636737}
    For $x,y,\delta$ in Proposition \ref{prop:6de7d130} we always have
    \begin{equation}\label{eq:841425e9}
        d^{\delta}_r(\lambda^ax)=\lambda^{a+r-n-1}y
    \end{equation}
    if $0\le a\le n$ and $0\le a+r-n-1<m-n$ (the differential is trivial if $a$ exceeds this range).
% END AIMPAPER STATEMENT
\uses{paper:4.6,paper:4.2}
% Blueprint note: Mathematical dependency review: [AimPaper source, lines 1138-1150](paper/main.tex.html#L1138). Multiply the boundary-extension formulas by lambda; the following remark identifies the complete classical differential coset.



% Correspondence: public counterparts; scope is specified below.
% Blueprint note: The public statements cover the full valid range of weights, the zero endpoint $a=n$, and vanishing when the source or target lies outside the strip.
% Blueprint note: This entry is bound to the existing Lean declarations. Its status reports their local completion; the scope of correspondence with the article is specified above.
% Blueprint note: Bound declarations (private names are resolved through stable source references without changing their visibility):
% Blueprint note: \begin{itemize}
% Blueprint note: \item \texttt{KIP126.ProofTools.Paper.syntheticDeltaExtension\_absolute}
% Blueprint note: \item \texttt{KIP126.ProofTools.Paper.syntheticDeltaExtension\_finite}
% Blueprint note: \item \texttt{KIP126.ProofTools.Paper.syntheticDeltaSource\_vanishes}
% Blueprint note: \item \texttt{KIP126.ProofTools.Paper.syntheticDeltaFiniteTarget\_vanishes}
% Blueprint note: \end{itemize}
\lean{KIP126.ProofTools.Paper.syntheticDeltaExtension_absolute,KIP126.ProofTools.Paper.syntheticDeltaExtension_finite,KIP126.ProofTools.Paper.syntheticDeltaSource_vanishes,KIP126.ProofTools.Paper.syntheticDeltaFiniteTarget_vanishes}

% Blueprint note: The original article number is 4.10; other numbering is generated by the blueprint. Throughout the article, spectra are assumed to be 2-completed, connective, and of finite type. Read this statement together with its original context.
% Blueprint note: \href{paper/main.tex.html#L1138}{Original statement, hypotheses, and context} · \href{paper/2412.10879.pdf}{Article PDF}
\end{corollary}
\begin{proof}
\textit{(See \cite{LWX126}, Corollary 4.10.)} \href{paper/main.tex.html#L1138}{Original statement and context}.

Apply multiplication by $\lambda^a$ to the boundary-extension formulas of Proposition~\ref{paper:4.6}. Naturality of the quotient cofiber sequence makes this multiplication commute with the connecting map; Proposition~\ref{paper:4.2} ensures that it introduces no additional positive-length extension. The source label is multiplied by $\lambda^a$ and the target label becomes $\lambda^{a+r-n-1}y$.

The two labels must lie in the source and target windows of Proposition~\ref{paper:3.12}. These are exactly the ranges displayed in the statement. Beyond those windows the relevant layer vanishes, so the corresponding extension is trivial. The quotient interpretation of the labels, including earlier boundaries, is the one fixed in Proposition~\ref{paper:4.6}.
\end{proof}

\section{Crossings of Adams differentials}

% AimPaper: 4.14 definition 1165 1173
\begin{definition}[Crossings of Adams differentials]
\label{paper:4.14}
\source{AimPaper Definition 4.14}
\level{coarse}
From \cite{LWX126}, Definition 4.14.\par
% BEGIN AIMPAPER STATEMENT
\label{def:classicalcrossdiff}
    Suppose $r\ge n+1$ and $d_r(x)=y$, where
    $$x\in E_2^{s,t}(X), \hspace{0.5cm} y\in E_2^{s+r,t+r-1}(X).$$
    A crossing of $d_r(x)=y$ on the $E_{n+1}$-page refers to an essential Adams differential
    $$d_{r-a-b}(x^\prime)=y^\prime,$$
    where
    $$x^\prime\in E_2^{s+a,t+a}(X), \hspace{.5cm} y^\prime\in E_2^{s+r-b,t+r-b-1}(X),$$
    with $0<a\le n-1$ and $0\le b\le r-n-1$. See Figure \ref{fig:f5eae42a}.
% END AIMPAPER STATEMENT
\uses{sem:SpectralSequence:differential,sem:SpectralSequence:cycles-boundaries}
% Blueprint note: Mathematical dependency review: [AimPaper source, lines 1165-1173](paper/main.tex.html#L1165). Define a classical crossing by a shorter essential Adams differential in the specified filtration rectangle.



% Correspondence: public counterparts; scope is specified below.
% Blueprint note: The definitions use the direction of classical differential crossings and the valid page range specified in the article. They are distinct from Moss crossings.
% Blueprint note: This entry is bound to the existing Lean declarations. Its status reports their local completion; the scope of correspondence with the article is specified above.
% Blueprint note: Bound declarations (private names are resolved through stable source references without changing their visibility):
% Blueprint note: \begin{itemize}
% Blueprint note: \item \texttt{KIP126.ProofTools.HasCrossingOn}
% Blueprint note: \item \texttt{KIP126.ProofTools.NoCrossingOn}
% Blueprint note: \end{itemize}
\lean{KIP126.ProofTools.HasCrossingOn,KIP126.ProofTools.NoCrossingOn}

% Blueprint note: The original article number is 4.14; other numbering is generated by the blueprint. Throughout the article, spectra are assumed to be 2-completed, connective, and of finite type. Read this statement together with its original context.
% Blueprint note: \href{paper/main.tex.html#L1165}{Original statement, hypotheses, and context} · \href{paper/2412.10879.pdf}{Article PDF}
\end{definition}

% AimPaper: 4.16 proposition 1218 1224
\begin{proposition}[Adams crossings and boundary-extension crossings]
\label{paper:4.16}
\source{AimPaper Proposition 4.16}
\level{coarse}
From \cite{LWX126}, Proposition 4.16.\par
% BEGIN AIMPAPER STATEMENT
\label{prop:cross-dr-En}
    The Adams differential $d_r(x)=y$ has a crossing on the $E_{n+1}$-page if and only if the corresponding $\delta_n$-extension
    $$d^{\delta_n}_r(x)=\lambda^{r-n-1}y$$
    for
    $$\delta_{n}:\nu X/\lambda^{n}\to \Sus^{1,-n}\nu X$$
    has a crossing.
% END AIMPAPER STATEMENT
\uses{paper:4.14,paper:2.7,paper:3.11,paper:3.12,paper:4.6,paper:4.10}
% Blueprint note: Mathematical dependency review: [AimPaper source, lines 1225-1234](paper/main.tex.html#L1225). Translate extension crossings through the infinity-page quotient coordinates and the lambda-multiple boundary formula.



% Correspondence: public counterparts; scope is specified below.
% Blueprint note: The public theorem gives the equivalence between classical crossings and crossings in the actual $\delta$-extension spectral sequence throughout the range $2\leq\mathrm{page}\leq r$.
% Blueprint note: This entry is bound to the existing Lean declarations. Its status reports their local completion; the scope of correspondence with the article is specified above.
% Blueprint note: Bound declarations (private names are resolved through stable source references without changing their visibility):
% Blueprint note: \begin{itemize}
% Blueprint note: \item \texttt{KIP126.ProofTools.Paper.classical\_crossing\_iff\_delta\_crossing}
% Blueprint note: \end{itemize}
\lean{KIP126.ProofTools.Paper.classical_crossing_iff_delta_crossing}

% Blueprint note: The original article number is 4.16; other numbering is generated by the blueprint. Throughout the article, spectra are assumed to be 2-completed, connective, and of finite type. Read this statement together with its original context.
% Blueprint note: \href{paper/main.tex.html#L1218}{Original statement, hypotheses, and context} · \href{paper/2412.10879.pdf}{Article PDF}
\end{proposition}
\begin{proof}
\textit{(See \cite{LWX126}, Proposition 4.16.)} \href{paper/main.tex.html#L1225}{Original proof}.

By Propositions \ref{prop:30e8b746}, \ref{prop:59f111f} and \ref{prop:6de7d130}, a crossing of $d^{\delta_n}_r(x)=\lambda^{r-n-1}y$ takes the form
    $$d^{\delta_n}_{r-a-b}(\lambda^a x')=\lambda^{r-n-1-b}y',$$
    where $0<a\le n-1$, $0\le b\le r-n-1$,
    $$\lambda^a x'\in E_\infty^{s+a,t+a,t}(\nu X/\lambda^n)\iso Z_{n-a}^{s+a,t+a}(X)/B_{1+a}^{s+a,t+a}(X),$$
    and
    $$y'\in E_\infty^{s+r-b, t+r-b-1}(\nu X).$$
    By Corollary \ref{cor:2a636737}, we see that this crossing corresponds to the classical Adams differential
    $$d_{r-a-b}(x')=y'.$$
\end{proof}

\section{Extension witnesses on finite and infinite pages}



% AimPaper: 5.9 definition 1437 1441
\begin{definition}[Crossings of $E_r$-extensions]
\label{paper:5.9}
\source{AimPaper Definition 5.9}
\level{coarse}
From \cite{LWX126}, Definition 5.9.\par
% BEGIN AIMPAPER STATEMENT
\label{def:41d51149}
    A crossing of the $(f,E_r)$-extension $d_{n}^{f,E_r}(x)=y$ in (\ref{eq:b98be76a}) is defined as an essential $(f, E_{r-a})$-extension from some $x'\in Z_{r-1-a}^{s+a,t+a}(X)$ to 
    $$y'\in Z_{r-1-n+b+e(f)}^{s+n-b,t+n-b}(Y)\backslash B_{1+n-b-e(f)}^{s+n-b,t+n-b}(Y)$$
    (where $\backslash$ denotes the difference of sets and it means that $y'$ should survive to the classical Adams $E_{r-n+b+e(f)}$ page while it should not be hit by an Adams differential of length at most $1+n-b-e(f)$) for $0<a \le r-2$ and $0\le b\le n-a-e(f)$. See Figure \ref{fig:66241bf6}.
% END AIMPAPER STATEMENT
\uses{paper:5.4}
% Blueprint note: Mathematical dependency review: [AimPaper source, lines 1437-1441](paper/main.tex.html#L1437). Define a crossing by an essential page extension with the displayed smaller-page and filtration bounds.



% Correspondence: public counterparts; scope is specified below.
% Blueprint note: Public definitions describe crossings and the absence of crossings for finite- and infinite-page extensions.
% Blueprint note: This entry is bound to the existing Lean declarations. Its status reports their local completion; the scope of correspondence with the article is specified above.
% Blueprint note: Bound declarations (private names are resolved through stable source references without changing their visibility):
% Blueprint note: \begin{itemize}
% Blueprint note: \item \texttt{KIP126.ProofTools.FiniteHasCrossing}
% Blueprint note: \item \texttt{KIP126.ProofTools.FiniteNoCrossing}
% Blueprint note: \item \texttt{KIP126.ProofTools.InfiniteHasCrossing}
% Blueprint note: \item \texttt{KIP126.ProofTools.InfiniteNoCrossing}
% Blueprint note: \end{itemize}
\lean{KIP126.ProofTools.FiniteHasCrossing,KIP126.ProofTools.FiniteNoCrossing,KIP126.ProofTools.InfiniteHasCrossing,KIP126.ProofTools.InfiniteNoCrossing}

% Blueprint note: The original article number is 5.9; other numbering is generated by the blueprint. Throughout the article, spectra are assumed to be 2-completed, connective, and of finite type. Read this statement together with its original context.
% Blueprint note: \href{paper/main.tex.html#L1437}{Original statement, hypotheses, and context} · \href{paper/2412.10879.pdf}{Article PDF}
\end{definition}

% AimPaper: 5.10 proposition 1515 1517
\begin{proposition}[Crossings of $E_r$-extensions and synthetic extensions]
\label{paper:5.10}
\source{AimPaper Proposition 5.10}
\level{coarse}
From \cite{LWX126}, Proposition 5.10.\par
% BEGIN AIMPAPER STATEMENT
\label{prop:cross-f-Er}
    An $(f,E_r)$-extension $d_{n}^{f,E_r}(x)=y$ (\ref{eq:b98be76a}) has a crossing if and only if the synthetic $\hat f_{r-1}$-extension $d_{n}^{\hat f_{r-1}}(x)=\lambda^{n-e(f)} y$ (\ref{eq:a020bad3}) has a crossing.
% END AIMPAPER STATEMENT
\uses{paper:5.9,paper:2.7,paper:3.12,paper:4.2,paper:2.15}
% Blueprint note: Mathematical dependency review: [AimPaper source, lines 1519-1536](paper/main.tex.html#L1519). Compare the two crossing definitions through multiplication by lambda, dividing its quotient-coordinate indeterminacy using the commutative square.

% Blueprint note: The available image-filtration and minimal-length lemmas are supporting steps, not the stated equivalence for arbitrary lengths. The direct formal binding below is a statement-only theorem with a registered pending proof; the local supporting results do not discharge it.


% Correspondence: original statement registered; proof pending.
% Blueprint note: Private proofs establish division by $\lambda$ in finite quotients of the same actual normalized map, an image-filtration estimate, and the absence of crossings in the minimal-length range. The equivalence between crossings at arbitrary lengths remains unproved. The available division theorem additionally requires target $\lambda$-kernel no-crossing and is specialized to the standard model; it does not prove this original equivalence with its fixed premises.


\sketch
% Blueprint note: The original article number is 5.10; other numbering is generated by the blueprint. Throughout the article, spectra are assumed to be 2-completed, connective, and of finite type. Read this statement together with its original context.
% Blueprint note: \href{paper/main.tex.html#L1515}{Original statement, hypotheses, and context} · \href{paper/2412.10879.pdf}{Article PDF}
% Blueprint note: Original-statement interface: \texttt{KIP126.ProofTools.Paper.finiteCrossing\_iff\_syntheticCrossing}. The statement is public; its proof and any required comparison remain pending.
\lean{KIP126.ProofTools.Paper.finiteCrossing_iff_syntheticCrossing}
\end{proposition}
\begin{proof}
\textit{(See \cite{LWX126}, Proposition 5.10.)} \href{paper/main.tex.html#L1519}{Original proof}.

\textit{Original argument and unresolved step.} The article writes a crossing in the large quotient in the form
\[
 d_{n-a-b}^{\widehat f_{r-1}}(\lambda^a x')
   =\lambda^{n-b-e(f)}y'.
\]
It then uses the naturality square for multiplication by $\lambda^a$ to try to divide this extension and obtain an extension in the quotient of length $r-1-a$. The resulting boundary ambiguity is compared with $B_{1+n-b-e(f)}$, and essentiality is intended to identify the two crossing conditions.

Naturality of the square does not by itself preserve the required image filtration after division: low-filtration $\lambda$-torsion can occur. Thus the printed argument is not a proof of the general equivalence as stated. The separate result at \ref{sem:ProofTools:finite-lambda-division} imposes the actual kernel no-crossing condition needed to preserve the same source and image. The issue and authorization boundary remain in the canonical AimPaper corrections record; the original comparison remains pending.
\end{proof}

% AimPaper: 5.14 proposition 1581 1587
\begin{proposition}[From infinity-page extensions to homotopy extensions]
\label{paper:5.14}
\source{AimPaper Proposition 5.14}
\level{coarse}
From \cite{LWX126}, Proposition 5.14.\par
% BEGIN AIMPAPER STATEMENT
Consider $f: X\to Y$, $x\in Z_\infty^{s,t}(X)$ and $y\in Z_\infty^{s+n,t+n}(Y)$.
    Then
    $$d^{f,E_\infty}_n(x)=y$$
    implies that
    $$d^{f}_n(x+B_\infty^{s,t}(X))=y+B_\infty^{s,t}(Y).$$
% END AIMPAPER STATEMENT
\uses{paper:5.4,sem:Literature:lambda-inversion,paper:3.11,paper:2.5,sem:Literature:realization-coordinates,sem:Literature:lambda-inversion-adams-comparison,sem:Literature:lambda-inversion-detection}
% Blueprint note: Mathematical dependency review: [AimPaper source, lines 1589-1597](paper/main.tex.html#L1589). Inverting lambda gives a map from the synthetic ESS to the classical ESS and sends the infinity-page extension to its homotopy extension.



As throughout \cite{LWX126}, spectra are understood to be 2-completed,
connective and of finite type. The Lean binding expresses the extension
equation by the representative criterion of Proposition~\ref{paper:2.5}:
$x$ detects one $a$ and $y$ detects $f_*a$. Realization coordinates and
the actual Adams-tower comparison of
Definition~\ref{sem:Literature:lambda-inversion-adams-comparison} make the
lambda-inversion background explicit. Zero labels, boundary labels and
inessential extensions remain allowed. Constructing these comparison
data for the fixed Standard model remains an upstream task, tracked in
issue 31.
% Blueprint note: The original article number is 5.14; other numbering is generated by the blueprint. Throughout the article, spectra are assumed to be 2-completed, connective, and of finite type. Read this statement together with its original context.
% Blueprint note: \href{paper/main.tex.html#L1581}{Original statement, hypotheses, and context} · \href{paper/2412.10879.pdf}{Article PDF}
% Blueprint note: The printed boundary group on $Y$ has index $(s,t)$ although $y$ has degree $(s+n,t+n)$. The typed conclusion uses the latter degree and keeps zero and inessential extensions. See the statement-typing correction record.
% Blueprint note: The AimPaper statement is retained verbatim and directly binds its proved representative-form transport. The realization coordinates and actual comparison remain explicit inputs; the binding does not assert that the Standard comparison producer has been constructed.
% Blueprint note: The deferred completed-background migration is tracked in [issue 31](https://github.com/SII-MATH/Kervaire126Main/issues/31).
\lean{KIP126.ProofTools.Paper.infiniteExtension_to_classicalExtension}
\end{proposition}
\begin{proof}
\uses{sem:ProofTools:extension-representatives,sem:Literature:lambda-inversion-detection,sem:Literature:lambda-inversion-adams-comparison,sem:AdamsSynthetic:normalized-map,paper:3.11,paper:2.5}
\textit{(See \cite{LWX126}, Proposition 5.14.)} \href{paper/main.tex.html#L1589}{Original proof}.

Invert $\lambda$ in the filtered synthetic map $\widehat f$. The comparison with the classical map $f$ induces a morphism of extension spectral sequences. On the starting page, Proposition~\ref{paper:3.11} sends a label in $Z_\infty/B_{1+t-w}$ to its class in $Z_\infty/B_\infty=E_\infty$.

For the formal representative transport, put $e=e(f)$ and $w=t+e$.
The actual extension equation gives one representative
$\beta\in\pi_{t-s,w}(\Sigma^{0,e}\nu X)$ whose image under $\hat f$
has the specified target detection. Apply the detection transport theorem
to this same $\beta$ and to $\hat f_*\beta$. The prescribed $E_2$
comparisons return precisely $x$ and $y$.

Realization of the normalized factorization, followed by cancellation of
the invertible realized lambda map, gives
\[
\theta_Y\circ\operatorname{Re}(\hat f)
 =f\circ\theta_{\Sigma^{0,e}\nu X}.
\]
Thus the first realized representative is $a$, and the second is exactly
$f_*a$. The sphere-coordinate transport accounts for
$(t+n)-(s+n)=t-s$. By Proposition~\ref{paper:2.5}, these detections give
the classical $f$-extension in the conclusion. The target is taken in its
own bidegree $(s+n,t+n)$, as required for the quotient by its boundary
group. Localization can enlarge the boundary indeterminacy, so the
resulting extension need not be essential. The argument includes zero
and boundary labels and imposes no nonvanishing or essentiality condition.
\end{proof}

\section{The finite-layer generalized Mahowald trick}

% AimPaper: 6.12 theorem 1780 1799
\begin{theorem}[The finite-layer generalized Mahowald trick]
\label{paper:6.12}
\source{A finite-layer revision of AimPaper, Theorem 6.12}
\level{coarse}
A finite-layer revision of \cite{LWX126}, Theorem 6.12.\par
The strong finite-layer assertion in Definition~\ref{sem:ProofTools:mahowald-law} holds for every compatible normalized triangle with exponent sum one. This includes both input levels, the actual no-crossing condition, the fixed strong quotient, and the same common representative and prescribed labels.
\uses{sem:ProofTools:mahowald-law,sem:ProofTools:mh-common-source,sem:Literature:normalized-finite-quotient-row,paper:6.11}
\sketch
\lean{KIP126.ProofTools.Standard.generalizedMahowaldAtLevel}
% Blueprint note: The complete finite-layer assertion is proved using the fixed Literature row, without a caller-supplied horizontal-row premise. It retains actual no-crossing synchronization, the same common representative, both signed image equations, both exact output detections, the original source cycle qualification and the specified target strong quotient. The proof introduces no local gap or axiom, does not use the frozen old MH or GL law, and still inherits the registered upstream model and comparison obligations.
% Blueprint note: This replaces the unprovable zero-layer version formerly displayed here. The original Theorem 6.12 remains in AimPaper and in the corrections page. Its frozen historical Lean law is not used as a proof of this revised target.
\end{theorem}
\begin{proof}
\textit{(See \cite{LWX126}, Theorem 6.12.)} \href{paper/main.tex.html#L1860}{Original Mahowald argument (before the finite-layer correction)}.

\textit{Proof of the finite-layer revision.} Fix the compatible normalized triangle and all data of Definition~\ref{sem:ProofTools:mahowald-law}. Actual no-crossing gives one representative $a$ for which the prescribed $h_A$ image and the $\delta_{A,B}$ image have the required detections. This synchronization is the common-source construction at \ref{sem:ProofTools:mh-common-source}; it does not replace either specified label.

Top-layer uniqueness identifies the boundary image of this $a$ with $g_B(b)$ for the chosen representative $b$. Apply the fixed finite-quotient horizontal row of \ref{sem:Literature:normalized-finite-quotient-row} and the triangle construction of Lemma~\ref{paper:6.11}. They produce a single $c$ satisfying both prescribed signed image equations. The comparison of the cofiber objects preserves both structure arrows and the same suspension comparison, so these are equations in the required quotients.

The normalized no-return argument removes the restriction defect of $c$. Visible-layer quotient reflection and the source/image detection comparisons then identify its two output detectors and retain the cycle qualification of the original source. Reading the second image equation in the specified strong quotient gives the finite extension in the conclusion. No $\lambda$ layer is cancelled, and the target is not projected to a smaller quotient. This proves the revised assertion relative to its registered upstream model and comparison inputs; it does not prove the zero-layer statement printed in the original Theorem 6.12.
\end{proof}

\section{Descent, indeterminacy and stretching}

% AimPaper: 6.20 proposition 1981 1993
\begin{proposition}[Restriction of extensions to earlier pages]
\label{paper:6.20}
\source{AimPaper Proposition 6.20}
\level{coarse}
From \cite{LWX126}, Proposition 6.20.\par
% BEGIN AIMPAPER STATEMENT
\label{prop:dec738d3}
    Suppose that we have an $(f,E_r)$-extension $d^{f,E_r}_n(x)=y$, where $x\in Z_{r-1}^{s,t}(X)$ and $y\in Z_{r-1-n+e(f)}^{s+n,t+n}(Y)$. Then for all $2\le r'<r$, we also have
    \begin{equation}\label{eq:dr'}
        d^{f,E_{r'}}_n(x)=y.
    \end{equation}

    Furthermore, if $d^{f,E_r}_n(x)=y$ is essential and $n\le r'-2+e(f)$, then (\ref{eq:dr'}) is inessential if and only if there exists some $0<a'\le n-e(f)$ and an element
    $$x'\in Z_{r'-1-a'}^{s+a',t+a'}(X)\backslash Z_{r-1-a'}^{s+a',t+a'}(X)$$
    such that
    $$d^{\hat f_{r'-1}}_{n-a'-b}(\lambda^{a'} x')=\lambda^{n-b-e(f)} y'$$
    for some $b\ge 0$.
% END AIMPAPER STATEMENT
\uses{paper:5.4,paper:2.18,paper:4.2,paper:3.12,paper:5.9}
% Blueprint note: Mathematical dependency review: [AimPaper source, lines 1995-2070](paper/main.tex.html#L1995). Restriction induces a morphism of ESSs; quotient coordinates and higher-source crossings describe the loss of essentiality.



% Correspondence: partial.
% Blueprint note: The corrected first assertion, including its required valid page-length bounds, is stated and bound separately in \href{#/Kervaire126#sem:ProofTools:finite-extension-descent}{Descent of finite-page extensions in the valid length range}. The printed second assertion, including its set-difference condition, remains unproved. A revised equivalence is proved publicly after removing that condition, fixing the original target $y$, and setting $b=0$; it retains the same actual earlier quotient, a classical source cycle, the original target, and all length bounds. This entry remains a partial correspondence with the original article.
% Blueprint note: This entry is bound to the existing Lean declarations. Its status reports their local completion; the scope of correspondence with the article is specified above.
% Blueprint note: The remaining related partial correspondence is \texttt{KIP126.ProofTools.Paper.finiteExtension\_inessential\_iff\_shorter}; the corrected descent theorem is bound only to its own mathematical statement.
\lean{KIP126.ProofTools.Paper.finiteExtension_inessential_iff_shorter}

\sketch
% Blueprint note: The original article number is 6.20; other numbering is generated by the blueprint. Throughout the article, spectra are assumed to be 2-completed, connective, and of finite type. Read this statement together with its original context.
% Blueprint note: \href{paper/main.tex.html#L1981}{Original statement, hypotheses, and context} · \href{paper/2412.10879.pdf}{Article PDF}
\end{proposition}
\begin{proof}
\textit{(See \cite{LWX126}, Proposition 6.20.)} \href{paper/main.tex.html#L1995}{Original proof}.

\textit{Restriction argument and scope of the correction.} Naturality of the quotient maps transports a representative extension in the length-$(r-1)$ quotient to the length-$(r'-1)$ quotient. The definition of an extension on $E_{r'}$ also requires its length to be legal there, namely $n\le r'-2+e(f)$. With this bound retained, the representative gives the descent theorem at \ref{sem:ProofTools:finite-extension-descent}. The unrestricted first assertion printed above omits this domain condition.

For indeterminacy, an inessential early-page extension has a higher-filtration source correction whose image accounts for its target. Expressing that correction in the finite-quotient cycle subquotients produces a shorter extension. This gives the separately stated corrected equivalence when the original target is retained and the set-difference requirement is removed. Requiring in addition that the shorter source fail to lift to the later page needs a further argument: a later-page crossing can obstruct that requirement.

The corrected dichotomy at \ref{sem:ProofTools:nonliftable-or-crossing} records these two alternatives. It does not establish the printed unrestricted equivalence with its general $b$, nor can the later-page extension's existence be assumed when proving a lifting theorem. These remaining boundaries are the ones recorded in the canonical corrections, not additional hypotheses silently inserted into the original statement.
\end{proof}

% AimPaper: 6.23 corollary 2074 2081
\begin{corollary}[Lifting extensions to later pages]
\label{paper:6.23}
\source{AimPaper Corollary 6.23}
\level{coarse}
From \cite{LWX126}, Corollary 6.23.\par
% BEGIN AIMPAPER STATEMENT
\label{cor:dfc6043e}
    Suppose $r\ge r'$ and there exists an $(f,E_{r'})$-extension 
    $$d^{f,E_{r'}}_n(x)=y$$
    for $x\in Z_{r-1}^{s,t}(X)$ and $y\in Z_{r'-1-n+e(f)}^{s+n,t+n}(Y)$. Assume that this extension has no crossing of the form $d^{f,E_{r'-a}}_{n-a-b}(x')=y'$ for any $b>0$, $0<a\le n-e(f)$, and 
    $$x'\in Z_{r'-1-a}^{s+a,t+a}(X)\backslash Z_{r-1-a}^{s+a,t+a}(X).$$ 
    Under these conditions, we also have an $(f,E_{r})$-extension
    $$d^{f,E_{r}}_n(x)=y.$$
% END AIMPAPER STATEMENT
\uses{paper:6.20,paper:5.9}
% Blueprint note: Mathematical dependency review: [AimPaper source, lines 2071-2081](paper/main.tex.html#L2071). Take the contrapositive of the restriction obstruction in Proposition 6.20 with the stated no-crossing condition.



% Correspondence: partial.
% Blueprint note: The current interface treats only a finite later page and requires the source and target to satisfy the later-page cycle conditions. It excludes every nonzero short extension in the same earlier finite quotient that does not lift, which is stronger than excluding only the essential crossings in the article. It does not establish the case $r=\infty$.
% Blueprint note: This entry is bound to the existing Lean declarations. Its status reports their local completion; the scope of correspondence with the article is specified above.
% Blueprint note: Direct formal counterparts: \texttt{KIP126.ProofTools.Standard.finitePageExtensionStretching}.
\lean{KIP126.ProofTools.Standard.finitePageExtensionStretching}

\sketch
% Blueprint note: The original article number is 6.23; other numbering is generated by the blueprint. Throughout the article, spectra are assumed to be 2-completed, connective, and of finite type. Read this statement together with its original context.
% Blueprint note: \href{paper/main.tex.html#L2074}{Original statement, hypotheses, and context} · \href{paper/2412.10879.pdf}{Article PDF}
\end{corollary}
\begin{proof}
\textit{(See \cite{LWX126}, Corollary 6.23.)} \href{paper/main.tex.html#L2074}{Original statement and context}.

\textit{The article's proposed deduction.} The text preceding this corollary identifies it as the contrapositive of the obstruction argument in Proposition~\ref{paper:6.20}: failure to retain an extension on the later page is supposed to produce one of the specified nonliftable short crossings. Excluding every such crossing would then force the lift.

This deduction requires the full obstruction statement, including legal extension lengths and control of later-page crossings; the corrected local results for Proposition~\ref{paper:6.20} do not alone supply it. The currently bound finite-page stretching theorem additionally requires the later-page cycle qualifications and excludes all relevant nonzero short obstructions, a stronger condition than the printed essential-crossing exclusion. It proves that qualified finite-page version, not the full printed corollary or the case $r=\infty$. Thus this block records the source argument and the precise remaining scope rather than claiming a completed proof of the original statement.
\end{proof}
